\quellseite{48}
\part[Part II. The Catechist in History]{Part II. Historical and Critical—The Catechist in History}

\quellseite{50}
\chapter[Chapter IV. The Catechist in Bible Times]{Chapter IV. The Catechist in History. In Bible Times}

Many Lutheran writers on Catechetics have professed to find a formal and complete catechumenate in the New Testament Church.\randnotiz{Formal catechumenate not in New Testament Church.} With them the wish is father to the thought, and the thought molds or invents the facts to suit. Nothing is ever gained by such procedure. A formal catechumenate did not exist in the New Testament Church. There is no trace of a regular, systematized course of instruction, to organized classes or to individuals previous to reception into the congregation by baptism. The Apostolic Church kept itself busy with missionary work. The Apostles went everywhere, preaching the Word. They preached mainly to the Jews who had Moses and the prophets. “For Moses of old time had in every city them that preach him, being read in the synagogues every Sabbath.” Jesus had ever “taught in the synagogue and in the temple, whither the Jews did always resort.” It was mainly to such synagogue-instructed Jews, familiar with the law and the prophets and the Psalms, or to Gentiles who had been influenced by the Jews and had become proselytes of the gate, or proselytes of righteousness, that the Apostles preached. Their hearers
\quellseite{51}
were not ignorant concerning the things of God. Most of them had come in contact with or had heard about John the Baptist and Jesus of Nazareth, and the first Apostolic and evangelistic preaching was to convince the hearers that Jesus was indeed the Messiah, the Christ, the Saviour of the world, in whom alone was remission of sins and salvation.

\randnotiz{Apostolic practice.} The facts as to Apostolic practice we find in the Book of Acts. After Peter’s Pentecost sermon, they that received the Word were baptized, and the same day there were added unto them about three thousand souls. As a result of Philip’s preaching in Samaria numbers were baptized, and there is no record of any systematic instruction or of a confession before baptism. To the eunuch Philip expounded a prophecy of Isaiah and applied it to Jesus. True, this was instruction, but it was not a course of catechising. Take the cases of Cornelius, the Jailer, Lydia, and others, and nowhere do we find an organized catechumenate with a systematic course of instruction. Further instruction, however, followed baptism, as in the case of Apollos (Acts 18:24, ff.). The congregations of believers were edified and strengthened by private and public study and instruction in the Old Testament Scriptures, in the unwritten traditions concerning Jesus and His redemptive work, and, later on, in the written accounts and in the Apostolic letters. The contrast
\quellseite{52}
between believers and unbelieving Jews and Gentiles became more and more marked. The lines of division were sharply drawn. Opposition, scorn and persecution became more bold and more bitter. The Christians were compelled to set forth their beliefs and teachings more clearly, to defend them against opposers and answer their objections. And so the necessity for clear, systematic teaching grew, and by the end of the first century there was a felt need of a Christian catechumenate.

But while we do not find an ordered catechumenate in operation in the Apostolic Church, the idea of religious nurture and instruction was there. The idea is older than the New Testament. Its roots lie far back in the Old Testament.\randnotiz{Religious nurture in Old Testament.}

paragraph

Inside the gates of Eden we find God Himself catechising His fallen children concerning commandment, sin, its consequences, its remedy. It was the earnest instruction of a yearning and loving Father in law and gospel, of sin and grace. And so God dealt directly not with Adam only, but with Cain, with Noah and with other antediluvians.

In the patriarchal age the father was the prophet, the teacher, the Catechist of his house, which often included several families with their servants. So in Gen. 18:19, God says of Abraham: “For I know him, he will command his children and his household after him that they shall keep the way of the Lord,
\quellseite{53}
to do justice and judgment.” This passage Matthew Henry explains thus: “Abraham not only took care of his children, but his whole household, including his servants, were catechised. In Ex. 12:26 ff., we have an example of catechisation on the meaning of the Passover. In Deut. 6:2 ff., the Lord, through Moses, clearly commands catechisation and gives the richest promises on its careful carrying out. (See also Josh. 4:6, 7, and 24:15; Ps. 34:11.)

In addition to this household instruction and catechisation, the Church of Israel also made provision for it. The priests and the Levites, in addition to their sacerdotal functions, were the official teachers or Catechists among the people. (See Lev. 10:11; Deut. 33:10; 2 Chron. 15:3; Ez. 44:23.)

In later times, when the Jewish nation became strong, outsiders were drawn towards the Jews. Before they were admitted to the privileges of the national religion they were instructed or catechised and afterwards, on confession of faith, received the so-called proselyte-baptism, and entered one or another of the different classes of proselytes.\footnote{See references and articles on Proselytes and on Synagogues in Edersheim’s “Life and Times of Jesus, the Messiah,” and in any good Bible Dictionary.}\randnotiz{Later Jewish instruction.}

In later days the synagogues also had regular instruction in the law and the prophets. There are writers, indeed, who claim that many synagogues in the principal cities had trained teachers from
\quellseite{54}
among the scribes, who had a daily afternoon service, free for all who would come, at which they would colloquially or catechetically expound their Scriptures.\footnote{Ibid.}

And so we see that while there is no formal modern catechumenate, with a Catechism, a class and a Catechist in the Old Testament, the idea is there. Religious instruction is as old as true religion, as old as humanity.

In the New Testament, as we have already seen, the formal catechumenate is also absent. But its seeds and germs are there. John the Baptist had been the nation’s teacher.\randnotiz{New Testament instruction.} Jesus Himself, according to the custom of His people, at twelve years old became a child of the law. We find Him in a kind of a catechetical exercise in the temple, both hearing His teachers and asking them questions. In the holy family in Nazareth He doubtless received daily instructions and impressions in the things of God. And so He increased in wisdom, and ripened up for his ministry. Then He went up and down in the land teaching on the highways, in the boat on the mountain side, instructing the multitude, His own disciples or the individual, in private home, in synagogue or in temple, everywhere endeavoring to make men wise unto salvation. See an example of His dealing with His own in Matt. 16: 13 ff.

\quellseite{55}
And while the Great Teacher was thus instructing all who would receive instruction, He was at the same time training the future teachers of the young Church. For a part of three years they were in His catechetical school.

Looking again into the teaching and practice of the Apostles we find the catechetical idea rooting and growing towards its future form.

Paul was a master in method. He was always tactful and methodical. He started with the known and made it lead into the unknown. He practiced \randnotiz{The apperceptive method of Paul.} what is called the modern pedagogical method of apperception. See examples of his tactful employment of the apperceptive method in Acts 14:14, ff., 17:16, ff., and 19:8, 9. The catechetical method of elementary instruction is more than hinted at, it is indeed taken for granted in the Epistles. (See 1 Cor. 3:1, 2; Heb. 5:11, 14; 6:1; 1 Pet. 2:2; 1 John 2:13.


\randnotiz{New Testament basis of the catechumenate.} The direct New Testament basis for the catechumenate we find in the great commission of Christ to the Church to make disciples (learners, followers) of all nations (Matt. 28:19, 20), and in Christ’s command to Peter, Feed my lambs (John 21:15).


Let us examine the great commission. Note the significant words by which it is framed in. The words, “All power is given unto me in heaven and in earth,” precede it. The impressive words, “And
\quellseite{56}
lo, I am with you alway, even unto the end of the world, Amen,” follow and close it. It is the all-powerful and ever-present Saviour who speaks in this His last great commission. Let us analyze it.

(a) Its great aim is μαθητεύσατε, i.e., make disciples.

(b) The means by which this is to be accomplished, βαπτίζοντες and διδάσκοντες, i.e., baptizing and teaching.

(c) The result of applying these means is seen in τηρεῖν πάντα i.e., those who have been made disciples observe all things whatsoever Christ has commanded.

It is the Lord’s will that this great commission shall embrace all nations, πάντα τὰ ἔθνη.

As we have seen above, the Apostles did not understand the commission as instituting a complete catechumenate. In their missionary labors among Jews and proselytes they were content to have their converts know and believe that Jesus is the Christ, the Saviour of the world. On this elementary knowledge of belief and confession they baptized them. After baptism they instructed them in the way of God more perfectly. The teaching and learning is, indeed, never finished, and never will be.

In the commission, then, we have the principles clearly laid down out of which the later formal catechumenate grew. The Apostles, on account of
\quellseite{57}
\randnotiz{How the formal catechumenate grew.} the nature of their converts, gained in succeeding ages, needed and used only an embryonic catechisation, preceding the baptizing of their adult converts. As the necessity grew the catechising increased in scope, in thoroughness and in rigor. But some elementary instruction always preceded baptism, which admitted to the fellowship and communion of the Church. The missionaries of necessity had to begin with adults; of these he that believed or professed to believe was baptized. But “How shall they believe in Him of whom they have not heard? And how shall they hear without one preaching?”

In the Apostolic age, then, the great commission was followed at first from its missionary aspect. That age had at most an embryonic missionary catechumenate, with adult catechumens, and, at first, adult baptisms. But the commission was intended for all nations. To the first adult converts it was said, “The promise is unto you and unto your children.” God Himself had instituted infant membership in His Church. God had never taken this back. It stood.

But infants are not born as disciples. They are not naturally born into the kingdom. They also must be made disciples.\footnote{See Chapters I. to V., “Way of Salvation in the Lutheran Church.”}
\quellseite{58}
\randnotiz{Children included.} The Apostles understood this. They also understood Mark 10:15, 16. Therefore, we soon find household baptisms in the Apostolic Church. These children are also to receive the baptizing and teaching which makes disciples. With them baptizing comes first and then teaching. From their baptism to their confirmation they ought to be in the Christian catechumenate. They are Christ’s lambs. The loving Shepherd commands the Church, “Feed my lambs.” The Church that catechises most carefully carries out this loving command most fully. The Lutheran Catechist has his feet planted solidly on the impregnable rock of the Holy Scriptures, and is in the Apostolic succession of Apostolic practice, which was in the line of the succession of patriarchs and prophets and Christ Himself. Truly a succession that the most mitred and titled prelate on earth might envy.


“A heritage, it seems to me,
A king might wish to hold in fee!”

\quellseite{59}
\chapter[Chapter V. The Catechist in History]{Chapter V. The Catechist in History. From the End of the Apostolic Age to the Reformation}

We come now to the Catechist in the post-Apostolic Church. In the Apostolic Church the services were free for all who would come.\randnotiz{The Catechist in the post-Apostolic Church.} The unbeliever, whether Jew or Gentile, had free access with the earnest inquirer and the believer. As opposition and persecution became more pronounced and bitter, the Christians felt themselves impelled to guard against hostile and lying spies, who would carry false reports of these services to the enemy. Hence the more and more careful guarding against dangerous intruders and the gradual growth of the \emph{Disciplina Arcani}. This fear of and guarding against dangerous spies also brought about more or less division and grading of the catechumens. In later times there has been much controversy as to this grading of the catechumens in the early Church. Some find four distinct classes,\randnotiz{Grading of catechumens.} viz.: (a) \emph{ἐξωθούμενοι}, i.e., those who were privately instructed and were not permitted to attend the public services of the congregation; (b) \emph{ἀκουόμενοι}, i.e., the hearers, who were admitted to hear the Scripture reading and sermon only; (c) \emph{γονυκλινοντες} i.e., those permitted to
\quellseite{60}
be present during the liturgical part of the service, to kneel during prayer and to join in the responses; (d) φωτιζόμενοι, or immediate candidates for baptism. This division is defended by Th. Harnack, Höfling, Zezschwitz, Kurtz, Krauszold, and Von Scheele. The division is disputed, at least it is claimed that it had no existence until after the time of Origen, or some time in the fourth century, by Mayer, Kliefoth, Hasselbach, Redepenning, Probst, Funk, Cohrs, Achelis, J. J. Rambach, and F. W. Schütze. Most of these admit a twofold division, i.e., that the catechumens, after they had been privately tested and had satisfied their teachers that they were sincere, had access to all the public services except to the celebration of the Lord’s Supper. To this holy mystery they were admitted as soon as they were baptized. After a somewhat detailed examination of the different views and arguments the writer of this at present accepts the latter view, admitting, however, that practices were by no means uniform. We have been impressed, especially by Achelis, who goes fully into this question and concludes that the writers who profess to find four classes of catechumens confound them with the penitentes, i. e., the lapsed ones who were seeking restoration to the Church and had to pass through four grades on their way back to the communion altar.\footnote{See “Lehrbuch der Practischen Theologie,” pp. 6-9.}
\quellseite{61}
\randnotiz{First mention of catechumenate.} As we study the catechumenate in the early Church, we find the first distinct mention in Justin Martyr, who died about the year 168, in his first great Apology, written about 140 A. D., during the persecution of Emperor Antonius Pius. In this important document Justin describes the worship of the congregation and the reception of members into its communion. We quote a few sentences:


“All who are convinced and believe that what we teach is true, and promise to live according to our teaching, are carefully instructed and directed to cry to God with prayer for the forgiveness of their sins, while we pray and fast with them.

“Then we lead them out to a place where there is water where they are born again even as we ourselves were.

“After the candidate has confessed his faith and has received baptism we lead him back into the congregation of our brethren, where earnest prayers are offered up for him, for ourselves and for our scattered brethren, that we who have received a knowledge of the truth might be diligent in good works and careful in observing God’s commandments, and so receive the salvation of our souls. After these prayers we greet one another with a holy kiss. Then the bread and wine are brought to the leader of the brethren, and a cup of wine is mixed with water. The leader offers praise and thanksgiving in the
\quellseite{62}
name of the Father, and of the Son, and of the Holy Ghost, who has made us meet to receive this heavenly gift. After this the deacons hand to each one present the bread and the wine and water which have been blessed. This feast we call the Eucharist. No one is permitted to partake of it but he who accepts what we teach and has the new birth in baptism and lives according to Christ’s commandment.”

\randnotiz{Facts derived from Justin Martyr.} From this interesting information, written less than fifty years after the death of the last of the Apostles, we learn these very important facts:

First. We see that the only division in the congregation was that between communicants and those who were under preparation for communicant membership.

Second. That the preparation consisted in instruction concerning the Father, Son and Holy Ghost, as it was probably set forth in a rudimentary form of what afterwards became the Apostles’ Creed. Whether the instruction was given by an official teacher or by a competent layman is not clearly stated.

Third. Justin does not inform us as to the length of the time of instruction. Clement of Rome, in the “Epistle to James,” which is probably spurious, is made to say that his own course lasted three months, but that some were admitted after a period of a few days only. It probably depended on
\quellseite{63}
the knowledge, character and spirit of the candidate. Our next important witness is \randnotiz{Tertullian’s witness.} Tertullian, who was a presbyter in Carthage, and died in the year 220.


In his time the catechumenate was further developed and systematized. It is he who, according to Zezschwitz, first uses the name “catechumenos.” Zezschwitz says: “Although we find the catechumenate in Tertullian’s time enriched on its liturgical side, it is essentially on the same basis as it was with Justin Martyr.”

Achelis says: “From all the accounts that we have it is clear that (in the time of Tertullian) the catechumens, from the time of their declaration of a desire to become Christians were counted in by the congregation and were objects of loving and prayerful attention.”\footnote{“Lehrbuch der Practischen Theologie,” p. 7.}

\randnotiz{Apostolic Constitutions’ witness.} From the Apostolic Constitutions, which probably date from the end of the third or the beginning of the fourth century, we learn that catechumens were instructed in the doctrine concerning God, His personality, His trinity, creation, providence, the nature and destiny of man, the redemptive agencies for the salvation of man and the final judgment. The more advanced were taught the mysteries of the incarnation and the person and work of the Saviour. The time of the catechumenate varied according to circumstances. The Apostolic Constitutions imply
\quellseite{64}
that it sometimes continued two and even three years. But they also declare that proper qualification and not length of instruction is to be the deciding factor.

Neither was there as yet a fixed rule as to who was to do the catechising. We find that this was done sometimes by the presbyter, sometimes by a deacon, and again by a lector or assistant of the presbyter. Redepening informs us in his book on Origen that deaconesses often instructed the women catechumens. As Christianity became more widely extended and attracted more and more attention, and as both inquirers and disputers came forward in ever larger numbers, it became more and more necessary to systematize and strengthen the catechumenate.


\randnotiz{Later growth.} In Alexandria a school for Catechists was established which demanded and obtained a line of teachers who were illustrious and far-famed for intellectual power, erudition and ability, as well as for spotless character and depth of piety. We need mention only the illustrious trio: Pantaenus, died about 190; Clement of Alexandria, died about 213, and Origen, died 254.

This renowned school, which at first instructed catechumens, became ere long a school for training Catechists and a theological seminary.\footnote{A most instructive 100-page pamphlet is “Die Katechetenschule zu Alexandria Kritisch Beleuchtet,” Von Fritz Lehmann, Leipzig, Alfred Lorentz, 1896.} In this
\quellseite{65}
school the weapons were forged with which to beat off the attacks of Greek Philosophy, Sophistry, Gnosticism, Neo-Platonism, Manichæism, and all sorts of heresy.\randnotiz{Training school for Catechists.} Here Greek met Greek, and the profoundest and acutest heathen opponents found foemen worthy of their steel. The teachers of this school and their better pupils could meet the opposer on his own ground and overcome him with his own weapons.

In this school the famous teachers instructed all who were willing to come. The method of teaching was by lecture, freely delivered, allowing question and discussion on the part of the hearer. Men and women crowded these lectures and discussions. The teachers received no stipulated salary, but lived from the free-will offerings of their pupils and patrons. Often they were very poor, lived the simple life, and were content with plain living and high thinking. The teacher would ask his pupils spontaneous questions and so awaken and prepare his mind to receive the truth he desired to impart and impress. The Bible was the great text-book. Free exposition and popular apology, earnest, direct application and ethical exhortation were freely intermingled. It required well-grounded, fearless and ready men to be good Catechists. For such men—men who at the same time had a burning zeal for
\quellseite{66}
Christ and His Church—such catechising was a most inspiring and blessed work.\footnote{We find other catechetical schools not so famous or efficient as the Alexandrian, in Ephesus, Smyrna, Antioch and Cæsarea.}

\randnotiz{Other citations.} We get a further glimpse of catechumens in Origen’s time from an extract of his great “Apology Against Celsus.” He says: “We Christians carefully prove those who desire to be admitted to our assembly, and only after they have given satisfactory evidence that they desire a virtuous and orderly life do we admit them after our usual custom. We form one class of beginners who have not yet received the symbol of renovation, and a second class of those who have proved their willingness to live as far as possible according to the teaching which we hold true.” The two classes were the catechumens and the communicants.

Cyril of Jerusalem, died 386, has left us twenty-three catechetical sermons. In the first eighteen he explains the general doctrines of Christianity, and in the last five the mysteries of the sacraments for those who were now ready to receive them.

Gregory of Nyssa, died 394, wrote a treatise for Catechists, in which he shows them how to meet and convince Jews, Gentiles and heretics of all classes and how to win the heart as well as the understanding. This brief survey covers the fourth century and its most important history.
\quellseite{67}
\randnotiz{The catechumenate after Constantine.} After the professed conversion of Constantine, when Christianity had become the State religion, hordes of heathen professed a desire to become Christians. A comparatively small proportion of these came with pure motives. It became necessary for the Church to exercise greater caution in the admission of members. The applicant for membership was closely questioned as to his motives. Inquiries were made as to his former life. If the reports were unsatisfactory the applicant was subjected to a severe testing before he was admitted to regular instruction. This, at least, was theory. But the state Church soon began to deteriorate, and carelessness in testing, in instruction and in admission became more and more frequent. The catechising was done acroamatically, as before, with opportunity for question and discussion.\footnote{See Augustin, “De Rudibus Catechizandis.”} The catechumen was instructed in the history of God’s kingdom, in the chief doctrine of Christianity and in its ethical demands. After the instruction, which was gradually lengthened and deepened, until it was a two or three years’ course, the candidate was asked whether he accepted the teaching, and if he answered affirmatively and was otherwise satisfactory, he received the sign of the cross and was admitted to immediate preparation for baptism.

One of the many false doctrines that began to
\quellseite{68}
creep in was the view that sins committed after baptism were specially heinous and so grave that the possibility of their forgiveness was at least doubtful, if not impossible.\footnote{Chrysostom calls out to a catechumen who was delaying: “How will you pledge yourself to lead a Christian life when life is at an end? You will be like a soldier who offers himself for service when the war is about to end.”} \randnotiz{False doctrine of baptism.} It was, therefore, deemed wise by many to postpone baptism until death seemed near. The emperor himself set a bad example in this direction. Besides the encouraging of the delay of baptism \randnotiz{Shortening of the catechumenate.} this false doctrine also had a dangerous tendency to increase carelessness and worldliness in the catechumens. Unsound doctrine begets an unsound life.

After the emperor had become nominally and officially Christian, and nearly every citizen wanted to be a professed Christian, there arose emphatic demand for the shortening of the catechumenate, and it was claimed that if one believed what the Church taught, or promised to believe the Church in the future, this ought to admit to baptism; that it was not necessary to insist on ethical requirements previous to baptism; that the proper time to teach ethics was after baptism. Unfortunately these voices gradually prevailed, and we have the sad spectacle of the decay of the catechumenate, the mass-baptism and the forced baptism of the middle ages. But even in the dark ages it was not all dark. There are
\quellseite{69}
bright gleams here and there amid the darkness. \randnotiz{Charlemagne’s efforts.} Charlemagne, as is well known, deplored the ignorance of the people and especially of the priests. He did all he could with the assistance of his cultured and brilliant prime minister, Alcuin, to plant Christian schools throughout his vast empire. In one of his public messages he says: “The priests should always admonish the people to learn the Symbol (i.e., the Apostles’ Creed), which is a seal of faith, and the Lord’s Prayer, and it is our will and command that a becoming discipline be exercised with those who refuse to learn, whether by fasting or by other penalties until they be willing to learn.” As we shall see, later efforts were not wanting to provide suitable text-books and catechisms for catechumens.

\randnotiz{Catechumenate for children.} Before we close this historical sketch it is important that we note how the missionary, adult catechumenate of the early Church, became an institution for children.

As we have already seen, infant baptism came into the early Church as naturally as it has come into every mission field not under the control of those who reject infant membership in the Christian Church.\footnote{For a condensed historical statement as to infant baptism in the Ancient Church, see Dr. Weidner’s article on Baptism in Luth. Cycl., and Whittekker’s valuable booklet, “Baptism.”}

The instruction of these children in the truth of
\quellseite{70}
God’s Word was left to the Christian parents. They were expected to bring them up in the nurture and admonition of the Lord. In the martyr Church this was conscientiously attended to. Beginning with the elect lady of the Second Epistle of John we meet with a number of illustrious Christian mothers, as, e.g., Emmelia, mother of Basil the Great; Anthusa, mother of Chrysostom; Monica, mother of Augustine.

When the Church became more and more worldly, cold and hierarchical, family life also deteriorated. But even in the darkest period Prince Georg of Anhalt could say: “The parents, especially the dear mothers, were often the best house-pastors and bishops.” When we come down to the Waldenses, the Wiclifites and the Hussites we find a beautiful Christian life in the family.

\randnotiz{Cloister instruction.} But back of all, we find the cloister in a manner substituting the family training. These monasteries and nunneries were not all and altogether bad. In many of them the light of learning and of piety was kept burning when all around was darkness. They were schools to which Christian parents brought their children for religious training. Basil the Great urged parents to send their children to these cloister schools. The Council of Mayence, 813, decreed that parents should send their children to the cloister schools. A later Synod at Mantua directed that a
\quellseite{71}
school should be erected beside every country church. But these, and other later laudable decrees, were generally dead letters, and were more conspicuous in the breach than in the observance. Still it was in the cloister and in the church schools that a kind of a child catechumenate was kept alive. The confessional was also a sort of a substitute, though a very poor one. The confessor was expected to instruct his children at least so much in the Ten Commandments, the Creed and the Lord’s Prayer that they might guide their confession by them.

\randnotiz{Lack of right religious instruction.} There was then a sad and sinful lack of the right kind of religious instruction. This was on the one hand a result of the worldliness of the hierarchy, and on the other hand a cause of the worldliness. The lack of a proper catechumenate made the terrible hierarchy possible and made the dark ages even darker until the dawn began to return. The decay of proper catechising made a reformation necessary.
